% =====================================================================
% PREAMBOLO PER TESI conforme alle NormeTesi
% =====================================================================

% Margini
\usepackage[a4paper,
    left=3.5cm,
    right=3cm,
    top=3cm,
    bottom=3cm]{geometry}

% Italiano
\usepackage[italian]{babel}
\usepackage[utf8]{inputenc}
\usepackage[T1]{fontenc}

% Interlinea
\usepackage{setspace}
\onehalfspacing % inter.: 1.5

% Header e Footer
\usepackage[headings]{fancyhdr}
\pagestyle{headings}

% Figure e grafici
\usepackage{graphicx}
\graphicspath{{immagini/}}
\usepackage{caption}
\usepackage{subcaption}
\usepackage{float}

% Tabelle
%\usepackage{tabularx}
%\usepackage{booktabs}

% Note a pié pagina
\usepackage{footmisc}
\renewcommand{\footnotesize}{\fontsize{10}{12}\selectfont}

% Citazioni lunghe (blockquote)
\usepackage{quoting}
\quotingsetup{
    vskip=6pt,
    leftmargin=20pt,
    rightmargin=20pt,
    font=small
}

% Matematica
\usepackage{amsmath}
\usepackage{amssymb}
\usepackage{amsfonts}
\numberwithin{equation}{chapter}

% Font Times + formula eleganti
\usepackage{newtxtext,newtxmath}

% Fisica
\usepackage{physics}
\usepackage{siunitx}
\usepackage[version=4]{mhchem} % reazioni 

% Microtipografia
\usepackage{microtype}

% Codici sorgenti
\usepackage{minted}
\usemintedstyle{friendly}
\setminted{
    linenos,        % numeri di riga
    breaklines,     % va a capo automaticamente
    fontsize=\small,
    frame=lines,    % cornice superiore e inferiore
    framesep=3mm,
    baselinestretch=1.1,
}

% Testo fittizzio
\usepackage{lipsum}

% Hyperref
\usepackage[hidelinks]{hyperref}

% Custom commands
